
On the minimal admissible branch the word-length candidates are
\begin{equation}
\label{eq:minimal-word-lengths}
\begin{gathered}
L_u^{\mathrm{diag}}=(L_u,L_c,L_t)=(7,3,0),\\
L_d^{\mathrm{diag}}=(L_d,L_s,L_b)=(7,5,2),\\
L_e^{\mathrm{diag}}=(L_e,L_\mu,L_\tau)=(8,5,3).
\end{gathered}
\end{equation}
These are not free parameters; they are the corrected minimal admissible diagonal word-length packets on the $C_6$ graph that keep the residual factors of \Cref{tab:yukawa-residuals} in the bounded $O(1)$ range. Once the hierarchy ordering is fixed, the admissible graph distance is determined up to full $6$-step windings by $\lambdaY\approx\num{0.2244}$.

\begin{remark}[Residue classes and full windings on the hexagon]
The exact resolvent formula \Cref{eq:finite-hexagonal-resolvent} makes the combinatorics behind \Cref{eq:minimal-word-lengths} transparent: every path length splits into a residue class on the base hexagon plus an optional full winding by $6$. Accordingly,
\[
L_t=0,\qquad
L_b=2,\qquad
L_\tau=L_c=3,\qquad
L_s=L_\mu=5
\]
live on the first hexagon, while
\[
\;L_u=L_d=1+6,
\qquad
L_e=2+6
\]
carry one additional full turn. Since
\[
\lambda_Y^6=\bigl(\phiz(1-\phiz)\bigr)^3\approx \num{1.28e-4},
\]
each extra winding produces the severe suppression needed for the first generation. This is not by itself a full mass theorem, but it is an exact algebraic reading of why the retained word-length pattern is so uneven.
\end{remark}

\begin{remark}[Universal $C_6$ transport triangle]
Modulo full windings, the three diagonal packets define the residue sets
\[
\{L_t,L_c,L_u\}\equiv\{0,3,1\},
\]
\[
\{L_b,L_s,L_d\}\equiv\{2,5,1\}=2-\{0,1,3\},
\qquad
\{L_\tau,L_\mu,L_e\}\equiv\{3,5,2\}=2+\{0,1,3\}
\pmod 6.
\]
Hence the up, down, and charged-lepton sectors lie on the same $D_6$ orbit of the base three-point subset $\{0,1,3\}\subset C_6$. In each case the cyclic distance multiset is exactly
\[
\{1,2,3\}.
\]
So the exact structure here is a three-point circular distance system with pairwise distinct cyclic distances, not a classical perfect difference set or a linear Golomb ruler.
Moreover, componentwise division by the full winding length gives
\[
\left\lfloor\frac{L_u^{\mathrm{diag}}}{6}\right\rfloor
=
\left\lfloor\frac{L_d^{\mathrm{diag}}}{6}\right\rfloor
=
\left\lfloor\frac{L_e^{\mathrm{diag}}}{6}\right\rfloor
=(1,0,0),
\]
so the first-generation direction carries exactly one extra full winding in every sector. This does not yet derive the diagonal packets dynamically, but it shows that any admissible replacement must preserve one universal $C_6$ transport triangle together with a common first-generation winding lock.
\end{remark}

\begin{principle}[Retained minimal-admissible complexity rule for diagonal transport packets]
\label{prin:minimal-admissible-transport-packets}
For each charge sector $f\in\{u,d,e\}$, let $\mathcal L_f$ denote the finite set of ordered triples
\[
L=(L_1,L_2,L_3)\in\mathbb N_{\mathrm{ord}}^3
\]
that satisfy strict hierarchy, positive-gap stability, the correct $C_6$ residue class of the sector, and bounded residual factors $\Lambda_{f,j}$. Then the retained diagonal packet is selected by the variational rule
\[
L_f^{\mathrm{diag}}
:=
\argmin_{L\in\mathcal L_f}\sum_{j=1}^3 L_j,
\]
and on the canonical branch the unique minimizers are
\[
L_u^{\mathrm{diag}}=(7,3,0),
\qquad
L_d^{\mathrm{diag}}=(7,5,2),
\qquad
L_e^{\mathrm{diag}}=(8,5,3).
\]
Thus the displayed word-length packets are retained branch data selected by the
minimal admissible complexity rule in the present draft, not yet a fully derived
dynamical theorem.
\end{principle}

\begin{proof}[Proof sketch]
The admissible search space is discrete and finite once hierarchy order, residue class, positive-gap stability, and bounded residual factors are imposed. Exhausting that finite list leaves one minimizer in each charge sector, and those minimizers are exactly the retained diagonal packets already displayed in \Cref{eq:minimal-word-lengths}. The universal $C_6$ transport triangle of the preceding remark is the rigid residue-level shadow of that enumeration.
\end{proof}

\begin{corollary}[Diagonal transport sum rule on the canonical branch]
\label{cor:transport-length-sum-rule}
Write
\[
L_f^{\mathrm{diag}}=(L_{f,1}^{\mathrm{diag}},L_{f,2}^{\mathrm{diag}},L_{f,3}^{\mathrm{diag}})
\qquad
(f\in\{u,d,e\}).
\]
Then on the canonical branch
\[
\sum_{f\in\{u,d,e\}}\sum_{j=1}^{3}L_{f,j}^{\mathrm{diag}}
=
\Oadm\gamma
=
10\,\bone-1
=
40.
\]
\end{corollary}

\begin{proof}
By \Cref{prin:minimal-admissible-transport-packets},
\[
\sum_{j=1}^3L_{u,j}^{\mathrm{diag}}=7+3+0=10,
\qquad
\sum_{j=1}^3L_{d,j}^{\mathrm{diag}}=7+5+2=14,
\]
\[
\sum_{j=1}^3L_{e,j}^{\mathrm{diag}}=8+5+3=16.
\]
Hence
\[
\sum_{f,j}L_{f,j}^{\mathrm{diag}}=10+14+16=40.
\]
On the other hand, \Cref{eq:41-chain-main} gives
\[
\Oadm\gamma=10\,\bone-1=40.
\]
Therefore the two quantities agree exactly.
\end{proof}

\begin{remark}[Status of the transport sum rule]
The equality of \Cref{cor:transport-length-sum-rule} is an exact cross-layer sum rule on the retained canonical branch. In the present manuscript it is not yet promoted to a conservation law, because no transport-dynamical theorem has been proved that would make the same sum invariant under an explicit time evolution or deformation flow.
\end{remark}

\begin{construction}[Least-curved representative of the family-holonomy class]
\label{cons:minimal-family-connection}
Let $\mathcal C_F$ be the space of $SU(3)_F$ connections on the family bundle
with fixed periods on the three fundamental cycles. Assuming a minimizer exists,
we denote by $\nabla_F$ a least-curved representative,
\[
\nabla_F \in \argmin_{\nabla\in\mathcal C_F}\|F_\nabla\|_{L^2}^2.
\]
In the present manuscript this is a retained variational gauge choice, not yet a
fully proved uniqueness theorem.
\end{construction}

\begin{remark}[Vacuum-overlap theorem path for the transport branch]
The retained transport grammar would become fully dynamical if the branch data were
generated by a global admissible vacuum functional
\[
\mathfrak V[g,\chigeo,\Phi,\nabla_F,\delta,\{L_f\}]
:=
\Gamma_{\mathrm{geom}}[g,\chigeo,\Phi]
+\mathcal F_{\mathrm{tr}}(\delta)
-\sum_f \log \mathcal O_f(L_f,\nabla_F)
+\Lambda_{\Seam}^{\mathrm{seam}},
\]
with vacuum overlaps
\[
\mathcal O_f(L,\nabla_F)
:=
\langle \Omega|
\Padm T_f(L,\nabla_F)^\dagger e^{-\beta H_{\mathrm{adm}}}
T_f(L,\nabla_F)\Padm
|\Omega\rangle.
\]
On the transport graph this overlap is expected to factor asymptotically as
\[
\mathcal O_f(L)\sim \lambdaY^{2L}\,|\Lambda_f(L)|^2,
\]
so the presently retained bounded residual factors become the statement that the
holonomy remainder stays in a physically admissible regime. Variation of
$\mathfrak V$ with respect to $\delta$, $\nabla_F$, $\Phi$, and the discrete
length data would then generate $\delta_{\mathrm{ph}}$, the least-curved family
connection, the Callias-selected Higgs branch, and the diagonal packets. In the
current manuscript this remains the explicit theorem path behind the retained
transport grammar.
\end{remark}

The transport-side status is summarized in \Cref{tab:transport-dictionary}.

\begin{table}[t]
\centering
\scriptsize
\resizebox{\linewidth}{!}{%
\begin{tabular}{@{}lllll@{}}
\toprule
Sector & Admissible length data & Residual prefactor & Role & Status \\
\midrule
charged leptons & $L_e^{\mathrm{diag}}=(8,5,3)$ & $\Lambda_e,\Lambda_\mu,\Lambda_\tau$ & ladder and hierarchy & bounded; closed in holonomy form \\
up quarks & $L_u^{\mathrm{diag}}=(7,3,0)$ & $\Lambda_u,\Lambda_c,\Lambda_t$ & up-type ratio ladder & bounded; closed in holonomy form \\
down quarks & $L_d^{\mathrm{diag}}=(7,5,2)$ & $\Lambda_d,\Lambda_s,\Lambda_b$ & down-type ratio ladder & bounded; closed in holonomy form \\
quark mixing & $L_{ij}^{q}$ & $\Lambda_{ij}^{q}$ & CKM transport entries & bounded; closed in holonomy form \\
neutrino mixing & $L_{ij}^{\nu}$ & $\Lambda_{ij}^{\nu}$ & PMNS transport entries & bounded; closed in holonomy form \\
hadron basis & center-neutral word lengths & block coefficients & meson and baryon transport basis & singlet closure on $\Hhad$ \\
strong CP & determinant spectrum of $\mathcal Y_f$ & none on the positive branch & phase suppression & closed on the admissible branch \\
\bottomrule
\end{tabular}}
\caption{Transport algebra dictionary with minimal admissible word-length candidates.}
\label{tab:transport-dictionary}
\end{table}

\subsection{CKM and PMNS sharpening}

The leading CKM magnitudes on the transport branch are
\begin{equation}
\label{eq:ckm-magnitudes}
s_{12}=\lambda_C,
\qquad
s_{23}=\frac{\phiz}{1+\lambda_C},
\qquad
s_{13}=\frac{\lambda_C^3}{3}.
\end{equation}
To second transport order, \Cref{prop:holonomy-triangle-phase} gives
\begin{equation}
\label{eq:ckm-phase}
\delta_{\mathrm{CKM}}=\frac{\pi}{3}+3\lambda_C^2+O(\lambda_C^4).
\end{equation}
The benchmark row used later keeps the truncation at that order,
\[
\delta_{\mathrm{CKM}}^{(2)}
:=
\frac{\pi}{3}+3\lambda_C^2
\approx
\num{1.198}\;\mathrm{rad},
\]
so $\pi/3$ remains the base monodromy phase while $3\lambda_C^2$ is the leading family-transport area correction.

For the neutrino sector the transport-side seesaw relation is kept as a geometric cross-check rather than the operational branch value used later:
\begin{equation}
\label{eq:neutrino-masses}
m_3^{\mathrm{ss}}=e^{2\lambda_C^2}\frac{\vgeo^2}{M_R},
\qquad
m_2^{\mathrm{ss}}=\pi\phiz\,m_3^{\mathrm{ss}},
\qquad
m_1^{\mathrm{ss}}=0,
\end{equation}
with $M_R\approx\num{1.31e15}\;\mathrm{GeV}$ fixed by the seesaw E$_8$ block at $(n,r,I_1)=(5,5,\frac{1}{12})$. Numerically this gives
\[
\!m_3^{\mathrm{ss}}\approx\num{4.939e-2}\;\mathrm{eV},
\qquad
\!m_2^{\mathrm{ss}}\approx\num{8.250e-3}\;\mathrm{eV},
\qquad
\Sigma m_\nu^{\mathrm{ss}}\approx\num{5.764e-2}\;\mathrm{eV}.
\]
\begin{remark}[Carrier-complement reading of the seesaw ratio]
On the retained seesaw cross-check branch,
\[
\frac{m_2^{\mathrm{ss}}}{m_3^{\mathrm{ss}}}
=
\pi\phiz
=
\pi\!\left(\frac{1-\gamma}{\pi}+\frac{3}{256\pi^4}\right)
=
(1-\gamma)+\frac{3}{256\pi^3}.
\]
Thus the leading hierarchy ratio is the carrier complement $1-\gamma=1/6$, corrected by a small seam term. In the present manuscript this is only a compact reading of the geometric seesaw cross-check; it is not promoted to an operational neutrino-branch theorem.
\end{remark}
The operational rows used later are kept separate: they use the oscillation-inversion normal-ordering branch
\[
m_1=0,
\qquad
m_2=\sqrt{\Delta m_{21}^2}\approx\num{8.614e-3}\;\mathrm{eV},
\qquad
m_3=\sqrt{\Delta m_{31}^2}\approx\num{5.015e-2}\;\mathrm{eV},
\]
so that
\[
\Sigma m_\nu=\num{5.8764e-2}\;\mathrm{eV}.
\]
This makes the bookkeeping explicit: oscillation inversion supplies the retained branch values, while the seesaw formula supplies a separate geometric cross-check through $M_R$.

\subsection{Confinement as admissible word selection}

\begin{proposition}[Admissible color words]
For each carrier word $w$, let $z(w)\in\{1,\omega,\omega^2\}$ be its color-center charge and let $\varepsilon(w)\in\{\pm1\}$ be its seam parity. Physical hadronic words satisfy
\[
z(w)=1,
\qquad
\varepsilon(w)=+1.
\]
\end{proposition}

The confinement claim is thus stated as an admissibility rule, not as a late dynamical slogan. The corresponding hadronic Hilbert space will be written explicitly in \Cref{sec:closure-theorems}. A schematic version is shown in \Cref{fig:admissible-words}.

\begin{figure}[t]
\centering
\begin{tikzpicture}[>=Latex, node distance=1.1cm and 1.3cm]
\node[draw=red!70!black, fill=red!4, rounded corners, minimum width=2.1cm, minimum height=1.0cm, align=center] (q) {free quark\\$z(w)\neq 1$};
\node[draw=green!50!black, fill=green!4, rounded corners, minimum width=2.1cm, minimum height=1.0cm, align=center, right=of q] (m) {meson\\$z(w)=1$};
\node[draw=green!50!black, fill=green!4, rounded corners, minimum width=2.1cm, minimum height=1.0cm, align=center, right=of m] (b) {baryon\\$z(w)=1$};
\draw[red!70!black, very thick] ($(q.north west)+(0.12,-0.12)$) -- ($(q.south east)+(-0.12,0.12)$);
\draw[red!70!black, very thick] ($(q.south west)+(0.12,0.12)$) -- ($(q.north east)+(-0.12,-0.12)$);
\node[below=0.2cm of q] {\small excluded};
\node[below=0.2cm of m] {\small admissible};
\node[below=0.2cm of b] {\small admissible};
\end{tikzpicture}
\caption{Confinement as admissible word selection.}
\label{fig:admissible-words}
\end{figure}

The transport-sector string tension is written as
\[
\sigma=\cthree^2\chi_{\QCD}^2,
\]
so the relative Wilson loop obeys
\[
\langle W(C)\rangle_{\mathrm{rel}}\sim e^{-\sigma A(C)}.
\]
The sector-dependent scale $\chi_{\QCD}$ is the minimal notation needed to avoid confusing the hadronic transport channel with the gravitational $\chigeozero$ background; the two channels are linked by the common underlying primitive response $\chiseed$ but not by a naive equality of scales.

\subsection{Transport-side ingredients for the strong-CP theorem}

Define
\begin{equation}
\label{eq:strong-cp}
\bar\theta=\theta_{\mathrm{sheet}}+\arg\det M_u+\arg\det M_d.
\end{equation}
\begin{definition}[Sheet-CP symmetry on the admissible branch]
\label{def:sheet-cp-symmetry}
Define the admissible sheet-CP operator by
\[
\mathcal C_\Sigma:=CP\circ\inv.
\]
\end{definition}

\begin{theorem}[Sheet-CP symmetry protection of strong CP]
\label{thm:sheet-cp-protection}
On the retained admissible branch, the relative action and topological charge satisfy
\[
\mathcal C_\Sigma\,S_{\mathrm{rel}}\,\mathcal C_\Sigma^{-1}=S_{\mathrm{rel}},
\qquad
\mathcal C_\Sigma\,Q\,\mathcal C_\Sigma^{-1}=-Q.
\]
Moreover, by \Cref{prop:gap-lift-full-operator} the regularized transport kernel induces a positive radial operator $H_f>0$ on the admissible sector, and the physical mass map admits the polar form
\[
M_f=U_{f,L}H_fU_{f,R}^\dagger,
\qquad
U_{f,L},U_{f,R}\in SU(3)_F.
\]
Hence
\[
\theta_{\mathrm{sheet}}=0,
\qquad
\arg\det M_f=0,
\]
while weak CP may still survive through
\[
V_{\mathrm{CKM}}=U_{u,L}^\dagger U_{d,L}.
\]
\end{theorem}

\begin{proof}[Proof sketch]
The involution $\inv$ implements the sheet exchange on the admissible branch, so composing it with $CP$ yields an exact symmetry of the relative action while flipping the sign of the topological charge. This forbids a bare $\theta$ deformation on the retained branch. The positive radial transport kernel removes intrinsic determinant phases, and the special-unitary holonomy factors preserve that vanishing while still allowing a nontrivial Jarlskog invariant in the relative family transport.
\end{proof}

\begin{corollary}[RG stability of $\bar\theta=0$]
\label{cor:strong-cp-rg-stability}
In every renormalization scheme that respects $\mathcal C_\Sigma$, $\Padm$, and the carrier structure, the vanishing of
\[
\bar\theta=\theta_{\mathrm{sheet}}+\arg\det M_u+\arg\det M_d
\]
is preserved under the declared RG and matching map.
\end{corollary}

The positive-branch clause is still essential here: the closure statement is about the physical mass map once the admissible transport branch has been fixed, not about an arbitrary unitary dressing of the same singular values. What changes in the present rewrite is the logical source of the null result: the branch no longer relies on a promise not to insert a bad counterterm by hand, but on the symmetry protection encoded by $\mathcal C_\Sigma$ together with the positive-kernel polar decomposition.

\begin{remark}[Admissibility as a unified three-sector motif]
The admissibility pair $(\mathsf{Adm},\Padm)$ acts as a single structural principle with three sector projections:
\begin{itemize}[leftmargin=1.5em, topsep=3pt, itemsep=2pt]
    \item \textbf{Electroweak:} seam-even selection picks the unique light Higgs doublet ($\inv\Phi=+\Phi$).
    \item \textbf{QCD:} center neutrality selects confined hadronic words ($z(w)=1$, $\varepsilon(w)=+1$).
    \item \textbf{Strong CP:} sheet-CP symmetry kills $\theta_{\mathrm{sheet}}$ in the relative action.
\end{itemize}
These are not three independent rules bolted onto different sectors. They are three projections of the same admissible-versus-inadmissible boundary: seam parity, center charge, and sheet exchange are all encoded in the double-cover and involution data of the primitive core. This unification is the conceptual backbone of the transport program.
\end{remark}

